\documentclass[../../../main]{subfiles}

\ifSubfilesClassLoaded{%
    \zexternaldocument*{../../1.Category-Theory/exported/aux/main}%
    \zexternaldocument*{../1.Additive-and-abelian-categories/exported/aux/main}%
    \zexternaldocument*{../2.Complexes-and-cohomology/exported/aux/main}%
    \zexternaldocument*{1.Mapping-cone}%
}{}

\begin{document}
\section{Quasi-isomorphisms}

\begin{definition}[Quasi-isomorphism] \label{def:quasi-isomorphism}
  A \hyperref[def:complex-morphism]{morphism of complexes} \(\alpha^{\bullet}\) is a \textbf{quasi-isomorphism} (abbreviated q.i.s) if it \hyperref[thm:LES-in-cohomology]{induces} an \hyperref[def:categorical-isomorphism]{isomorphism} in \hyperref[def:cohomology]{cohomology}.
\end{definition}

\begin{remark}
  By \autoref{cor:exactness-mapping-cone}, \(\alpha^{\bullet}\) is a \hyperref[def:quasi-isomorphism]{quasi-isomorphism} if and only it its \hyperref[def:mapping-cone]{mapping cone} \(\mCone{\alpha}{\bullet}\) is an \hyperref[def:exact-complex]{exact complex}.
\end{remark}

\begin{example}[Quasi-isomorphisms are not isomorphisms]
  In \(\Complex{\Ab}\), consider the \hyperref[def:complex-morphism]{morphism of complexes}
  \begin{center}
    \begin{tikzcd}
      \cdots \arrow[r] & 0 \arrow[r] \arrow[d] & \mathbb{Z} \arrow[r, "\cdot2"] \arrow[d] & \mathbb{Z} \arrow[r] \arrow[d, "\pi"] & 0 \arrow[r] \arrow[d] & \cdots \\
      \cdots \arrow[r] & 0 \arrow[r] & 0 \arrow[r] & \quot{\mathbb{Z}}{2\mathbb{Z}} \arrow[r] & 0 \arrow[r] & \cdots
    \end{tikzcd}
  \end{center}
  where \(\pi\) is the obvious projection.
  Note that this is a \hyperref[def:quasi-isomorphism]{quasi-isomorphism}.
  However, it is not an \hyperref[def:categorical-isomorphism]{isomorphism} because there are no (non-trivial) group homomorphisms \(\quot{\mathbb{Z}}{2\mathbb{Z}} \to \mathbb{Z}\).
\end{example}

\begin{lemma}
  Let \(\category{A}\) be an \hyperref[def:additive-category]{additive category} and let \(F: \Complex{\category{A}} \to \category{D}\) be an \hyperref[def:additive-functor]{additive functor} such that, for every \hyperref[def:quasi-isomorphism]{quasi-isomorphism} \(\rho^{\bullet}\), \(F(\rho^{\bullet})\) is an \hyperref[def:categorical-isomorphism]{isomorphism}. Let \(M^{\bullet}\) be an \hyperref[def:exact-complex]{exact complex} in \(\Complex{\category{A}}\).
  \begin{itemize}
    \item The \hyperref[def:complex]{complex} \(F(M^{\bullet})\) is a \hyperref[def:zero-object]{zero object} in \(\category{D}\).
    \item Let \(\alpha^{\bullet}: L^{\bullet} \to N^{\bullet}\) be a \hyperref[def:complex-morphism]{morphism of complexes} in \(\Complex{\category{A}}\) that factors through an \hyperref[def:exact-complex]{exact complex}
          \begin{center}
            \begin{tikzcd}[column sep=tiny]
              L^{\bullet} \arrow[rr, "\alpha^{\bullet}"] \arrow[dr] & & N^{\bullet} \\
              & M^{\bullet} \arrow[ur]
            \end{tikzcd}
          \end{center}
          Then \(F(\alpha^{\bullet})\) is the zero morphism.
  \end{itemize}
\end{lemma}
\begin{proof}
  Since \(M^{\bullet}\) is \hyperref[def:exact-complex]{exact}, the zero morphism \(M^{\bullet} \to M^{\bullet}\) is a \hyperref[def:quasi-isomorphism]{quasi-isomorphism}.
  Therefore we have
  \begin{center}
    \begin{tikzcd}
      F(M^{\bullet}) \arrow[r, "F(0)"] \arrow[rr, bend right, "\id_{F(M^{\bullet})}"'] & F(M^{\bullet}) \arrow[r, "F(0)^{-1}"] & F(M^{\bullet})
    \end{tikzcd}
  \end{center}
  Since \(F\) is \hyperref[def:additive-functor]{additive} \(F(0) = 0\), making \(\id_{F(M^{\bullet})} = 0\).

  For the other part, applying \(F\) to the diagram we get
  \begin{center}
    \begin{tikzcd}[column sep=tiny]
      F(L^{\bullet}) \arrow[rr, "F(\alpha^{\bullet})"] \arrow[dr] & & F(N^{\bullet}) \\
      & F(M^{\bullet}) \arrow[ur]
    \end{tikzcd}
  \end{center}
  Therefore since \(F(M^{\bullet})\) is a \hyperref[def:zero-object]{zero object} in \(\category{D}\) we have that \(F(\alpha^{\bullet}) = 0\).
\end{proof}

\end{document}
